\title{The Era of 1-bit LLMs: All Large Language Models are in 1.58 Bits}

\begin{abstract}
Building on advancements exemplified by BitNet~\cite{bitnet}, the landscape of Large Language Models (LLMs) is transitioning towards 1-bit representations. In this paper, we present a novel 1-bit LLM, called \textbf{\text{BitNet b1.58}}, where every model parameter takes a ternary value from the set \{-1, 0, 1\}. This architecture is able to achieve performance parity—measured both in perplexity and task outcomes—with traditional full-precision (FP16 or BF16) Transformer LLMs, when matched for parameter count and training data. Crucially, this ternary model dramatically outperforms its counterparts in terms of computational cost, reducing requirements for memory, latency, throughput, and energy use. On a broader scale, the introduction of the 1.58-bit LLM establishes a revised scaling principle and an updated training framework that facilitates the creation of future LLMs which excel in efficiency and performance. Moreover, it fosters novel computational models and lays a foundation for hardware architectures specifically designed for efficient 1-bit LLM operation.
\end{abstract}

\section{The Era of 1-bit LLMs}

The recent trajectory of artificial intelligence research has been characterized by a dramatic expansion in the scale and functionality of Large Language Models (LLMs). While these models have achieved exceptional results across various NLP benchmarks, their rapidly growing sizes introduce significant hurdles to practical deployment, accompanied by substantial concerns regarding the financial and environmental costs associated with their high energy demands.

One solution to mitigate these issues involves leveraging post-training quantization to yield reduced-bit models suitable for inference~\cite{smoothquant, gptq, quip, quip_sharp}. This method lowers the arithmetic precision for model weights and activations, offering substantial gains in computational and memory efficiency. The prevailing movement has shifted from standard 16-bit to more compact 4-bit designs~\cite{gptq, awq}. However, although post-training quantization remains a common approach in deploying LLMs across industry, it is inherently limited and does not deliver optimal results.

Emerging research into 1-bit model designs, such as BitNet~\cite{bitnet}, highlights a compelling path toward lowering the operational expense of LLMs while preserving their effectiveness. Traditional LLMs utilize 16-bit floating-point representations (FP16 or BF16), with matrix multiplication comprising the dominant share of computation. This results in substantial computational overhead primarily due to floating-point addition and multiplication. BitNet, by contrast, executes matrix multiplications using only integer addition, yielding significant reductions in energy consumption for LLMs. Since power serves as a limiting factor for performance in numerous hardware architectures, this energy efficiency readily translates into enhanced computational speed.

Besides the arithmetic computations, another significant bottleneck is the movement of model parameters from DRAM to on-chip memory resources like SRAM during inference. Expanding SRAM to accelerate throughput has been explored, though it incurs much greater cost compared to DRAM. Relative to full-precision counterparts, 1-bit LLMs substantially decrease the memory requirements, both in terms of storage capacity and data bandwidth, thus reducing both the expense and the latency associated with weight loading from DRAM, ultimately enabling faster, more efficient inference.

In this study, we propose a notable 1-bit LLM extension called \textbf{\text{BitNet b1.58}}. In this approach, every model parameter is allowed three possible values: \{-1, 0, 1\}. By incorporating a 0 into the binary space of the original 1-bit BitNet, the representation effectively increases to 1.58 bits. \text{BitNet b1.58} preserves the full spectrum of benefits seen with BitNet, such as the novel computation scheme that eliminates nearly all multiplications during matrix operations and is amenable to extensive hardware optimizations. Energy consumption remains on par with standard BitNet, and it demonstrates much improved efficiency over FP16 LLMs in terms of memory usage, processing throughput, and latency. Furthermore, \text{BitNet b1.58} introduces two additional strengths: First, its explicit use of 0 in the weight values enables feature selection, thereby enhancing the expressive power of the model and improving the practical effectiveness of 1-bit LLMs. Second, empirical evaluations indicate that starting from a 3B parameter scale, \text{BitNet b1.58} can reach parity with full-precision (FP16) LLMs, matching their perplexity and end-task outcomes given identical training setups (e.g., model scale, amount of training data, etc.).

\section{BitNet b1.58}

\text{BitNet b1.58} is derived from the BitNet framework, a variant of the Transformer architecture that substitutes \emph{nn.Linear} layers with \emph{BitLinear}. This model is trained from the beginning using weights quantized to 1.58 bits and activations quantized to 8 bits. Relative to the standard BitNet, several changes are implemented, outlined as follows.

\paragraph{Quantization Function.} To ensure the weight values are limited to -1, 0, or +1, an \emph{absmean} quantization function is employed. This method rescales the weight matrix based on its mean absolute value, subsequently rounding each element to the nearest integer from the set \{-1, 0, +1\}:

\begin{equation}
    \widetilde{W} = \mathrm{RoundClip}(\frac{W}{\gamma + \epsilon}, -1, 1),
\end{equation}

\begin{equation}
    \mathrm{RoundClip}(x, a, b) = \max(a, \min(b, \mathrm{round}(x))),
\end{equation}

\begin{equation}
    \gamma = \frac{1}{nm}\sum_{ij} |W_{ij}|.
\end{equation}

The activation quantization procedure remains consistent with that used in BitNet, with the distinction that activations are not scaled to the interval $[0, Q_b]$ prior to applying non-linearities. Instead, activations for each token are rescaled to the interval $[-Q_b, Q_b]$, effectively eliminating zero-point quantization. This approach simplifies both implementation and system-level optimization, and, in our empirical evaluations, any impact on model performance is negligible.

\paragraph{LLaMA-alike Components.}

LLaMA~\cite{llama, llama2} has become the foundational architecture for most open-source LLMs. To foster compatibility with open-source initiatives, \text{BitNet b1.58} incorporates elements from LLaMA. Concretely, it utilizes RMSNorm~\cite{rmsnorm}, SwiGLU~\cite{swiglu}, rotary embedding~\cite{rope}, and omits bias terms entirely. This design choice ensures that \text{BitNet b1.58} can be readily adopted within widely used open-source platforms (such as Huggingface, vLLM~\cite{vllm}, and llama.cpp\footnote{\url{https://github.com/ggerganov/llama.cpp}}) with minimal adaptation required.

\section{Results}

We evaluated \text{BitNet b1.58} against our reproduced FP16 \text{LLaMA LLM} models across multiple scales. For consistency and fair assessment, all models underwent pre-training on the RedPajama dataset~\cite{redpajama}, using 100 billion tokens. Zero-shot capabilities were assessed on a suite of language benchmarks: ARC-Easy~\cite{arc}, ARC-Challenge~\cite{arc}, Hellaswag~\cite{hellaswag}, Winogrande~\cite{winoGrande}, PIQA~\cite{piqa}, OpenbookQA~\cite{openbookqa}, and BoolQ~\cite{boolq}. We additionally reported validation perplexity using the WikiText2~\cite{wikitext2} and C4~\cite{c4} datasets.

We further conducted empirical comparisons of runtime GPU memory consumption and inference latency between \text{BitNet b1.58} and \text{LLaMA LLM}. Measurements were collected via the FasterTransformer\footnote{\url{https://github.com/NVIDIA/FasterTransformer}} framework, recognized for its efficiency in LLM inference on GPU hardware. For \text{BitNet b1.58}, the 2-bit kernel from Ladder~\cite{ladder} was utilized. Our primary inference metric was the generation time per output token, representing the principal cost during inference.

As presented in Table~\ref{tab:ppl}, we summarize both perplexity results and computational resource requirements for \text{BitNet b1.58} in comparison to \text{LLaMA LLM}. Notably, beginning at the 3B parameter scale, \text{BitNet b1.58} achieves perplexity on par with the full-precision \text{LLaMA LLM}, while delivering a 2.71$\times$ reduction in generation latency and a 3.55$\times$ decrease in GPU memory usage. Specifically, the 3.9B variant of \text{BitNet b1.58} not only surpasses the performance of \text{LLaMA LLM} 3B but does so with 2.4$\times$ the speed and 3.32$\times$ less memory utilization.

Detailed zero-shot task performance is presented in Table~\ref{tab:zero-shot}, with evaluations conducted following the procedure outlined by \emph{lm-evaluation-harness}\footnote{\url{https://github.com/EleutherAI/lm-evaluation-harness}}. The findings indicate that as model scale increases, the performance discrepancy between \text{BitNet b1.58} and \text{LLaMA LLM} diminishes. Crucially, starting at the 3B parameter scale, \text{BitNet b1.58} matches the accuracy of the full-precision baseline. Consistent with perplexity observations, \text{BitNet b1.58} 3.9B demonstrates superior performance to \text{LLaMA LLM} 3B while incurring substantially lower latency and memory cost. These results collectively position \text{BitNet b1.58} as a Pareto improvement relative to leading LLM architectures.

\paragraph{Memory and Latency}

The model was also extended to larger sizes of 7B, 13B, and 70B parameters to assess its associated computational cost. As depicted in Figure~\ref{fig:latency-memory-energy}, both latency and memory requirements exhibit consistent trends with scaling, where performance gains, specifically speed-ups, become more pronounced for larger architectures. Notably, \text{BitNet b1.58} at 70B parameters outpaces the \text{LLaMA LLM} reference by a factor of 4.1 in terms of speed. This improvement is attributed to the increasing computational overhead of the \emph{nn.Linear} operation as models grow larger. Memory usage behaves analogously, with the embedding layer’s contribution to overall memory decreasing for larger models since it is always stored in full precision. The reported latency and memory usage were obtained using a 2-bit kernel, indicating potential for further optimizations and cost reductions.

\paragraph{Energy}

An evaluation of arithmetic operation energy requirements was also conducted for both \text{BitNet b1.58} and \text{LLaMA LLM}. The focus was primarily on matrix multiplication, which constitutes the primary computational expense in LLMs. Figure~\ref{fig:energy} details the breakdown of these energy expenditures, demonstrating that the majority of \text{BitNet b1.58}'s consumption derives from INT8 addition, while \text{LLaMA LLM} relies on both FP16 addition and multiplication. Utilizing the energy modeling approaches from~\cite{energycost, pokebnn}, it was determined that \text{BitNet b1.58} requires 71.4 times less energy for matrix multiplication compared to its FP16-based counterpart when executed on 7nm hardware. Additional reporting includes end-to-end energy use for processing sequences of 512 tokens. The data indicate that as model size increases, \text{BitNet b1.58} demonstrates increasingly superior energy efficiency over the FP16 \text{LLaMA LLM} baseline. This is primarily due to the expanding contribution of \emph{nn.Linear} operations to total cost at scale, with other modules accounting for a diminished share in larger models.

\paragraph{Throughput}

To evaluate throughput, we benchmarked \text{BitNet b1.58} and \text{LLaMA LLM} with 70B parameters on two 80GB A100 GPUs, utilizing pipeline parallelism~\cite{gpipe} to facilitate execution of the \text{LLaMA LLM} 70B model on the hardware. The batch size was systematically increased up to the limit permitted by GPU memory, using a fixed sequence length of 512. As indicated in Table~\ref{tab:throughput}, \text{BitNet b1.58} 70B supports up to eleven times the batch size of \text{LLaMA LLM}, leading to an overall throughput increase by a factor of 8.9.

\textbf{\text{BitNet b1.58} introduces a new paradigm in scaling laws, particularly in the tradeoff between model efficacy and inference efficiency}. Based on the findings shown in Figure~\ref{fig:latency-memory-energy} and \ref{fig:energy}, the following correspondences can be drawn between 1.58-bit and 16-bit models in terms of their respective sizes:

\begin{itemize}
    \item The 13B parameter \text{BitNet b1.58} outperforms a 3B FP16 LLM regarding latency, memory footprint, and power consumption.
    \item The 30B \text{BitNet b1.58} offers greater efficiency—in terms of latency, memory, and energy—than a 7B FP16 LLM.
    \item The 70B \text{BitNet b1.58} is likewise more efficient with respect to latency, memory usage, and energy needs compared to a 13B FP16 LLM.
\end{itemize}

\paragraph{Training with 2T Tokens}

The volume of training tokens significantly impacts LLM performance. To examine the token scaling properties of \text{BitNet b1.58}, we trained a model using 2T tokens, adopting the data pipeline described in StableLM-3B~\cite{StableLM-3B-4E1T}, a leading open-source 3B parameter model. Both \text{BitNet b1.58} and StableLM-3B were assessed using a suite of tasks: Winogrande~\cite{winoGrande}, PIQA~\cite{piqa}, SciQ~\cite{sciq}, LAMBADA~\cite{lambada}, and ARC-easy~\cite{arc}. Zero-shot accuracy metrics are provided in Table~\ref{tab:2t}. Where both accuracy and normalized accuracy were available, we report their mean. The results for StableLM-3B with 2T tokens are extracted directly from its technical documentation. The data demonstrate that \text{BitNet b1.58} consistently surpasses StableLM-3B across all benchmark tasks, suggesting that LLMs utilizing 1.58-bit quantization not only excel in efficiency but also retain robust generalization capabilities.

\section{Discussion and Future Work}

\textbf{1-bit Mixture-of-Experts (MoE) LLMs}

Mixture-of-Experts (MoE) models are widely recognized for their efficiency in large language models (LLMs). Although they effectively decrease computational FLOPs, their extensive memory demands and substantial inter-chip communication costs present obstacles for large-scale deployment. These limitations can be mitigated by using 1.58-bit LLMs. By minimizing the memory requirements, fewer hardware devices are necessary to serve MoE-based models. Furthermore, the reduced model size lowers the data transfer overhead across devices. Ideally, should the entire model fit on a single chip, interconnect overhead could be eliminated completely.

\textbf{Native Support of Long Sequence in LLMs}

The capability to process long sequences has become a key requirement for contemporary LLMs. A primary bottleneck for inference with lengthy sequences is the memory overhead incurred by storing KV caches. The introduction of \text{BitNet b1.58} marks an important advancement by cutting the activation bitwidth from 16 to 8 bits, thereby permitting a doubling of the maximum context length under equivalent resource constraints. Further, there is potential to compress these representations losslessly down to 4 bits or even lower in 1.58-bit LLMs, a direction we propose for future research.

\textbf{LLMs on Edge and Mobile}

Adopting 1.58-bit LLMs could drastically boost the efficiency of language models on edge and mobile hardware, where memory and compute limitations traditionally hinder LLM deployment. The lower memory footprint and decreased power consumption of 1.58-bit LLMs make it feasible to bring advanced language capabilities to these platforms, unlocking new applications previously out of reach. This extension of LLM functionalities to edge and mobile opens up novel user experiences. In addition, 1.58-bit LLMs are more compatible with CPU-based environments—prevalent in edge and mobile contexts—so \text{BitNet b1.58} can run efficiently on such devices, enhancing their functionality further.

\textbf{New Hardware for 1-bit LLMs}

Emerging platforms such as~Groq\footnote{\url{https://groq.com/}} have shown the efficacy of specialized architectures, like LPUs, tailored for LLM computation. Building upon this trend, we advocate for the exploration and development of dedicated hardware and systems optimized for 1-bit LLMs, capitalizing on the new computation schemes made possible by BitNet~\cite{bitnet}.